% !TeX root = main.tex

\chapter{三月七日}
  京都、JR山科駅前。昼飯に最適なご飯屋というと、踏切を渡って、左手に位置する松屋、あるいはその先の大阪王将である。

  さて、今日はどちらにしようかとまずは松屋を目的地に設定し、歩き出す。踏切をわたり終えた後、何気なく腕時計を見る。一時三十分、次の用事まで後三十分であった。これはいけないと、目的地設定をいったん解除し、ぶらぶらと歩きながら次の手を考える。

  あれこれ考えた末、禁じ手を使う。みんな大好きコンビニエンスストアである。短時間で昼食を済ませたいときはコンビニを頼ってしまう。

  店内に入ると、アジア系の外国人女性の店員がにこやかに会釈をした。コンビニに入るたびに、この国は外国人の方々に支えられているなと感じると同時に、自分はいつまで日本で引きこもりをやっているのだと少し嫌な気分になる。
  そんなことを考えながら、サンドイッチと温かいお茶を手に取り、タッチ決済で迅速に会計を済ます。

  店を出るとすぐ近くに良い具合の公園があった。日差しの当たっているベンチは山登り前のご老人方に占領されていたので、しかたなく日陰で冷たくなったベンチに腰掛ける。

  レタスとハムの入った、多少健康志向のサンドイッチをほおばる。コンビニ店員だったときには期限切れのサンドイッチばかり食べていたので、久々の期限内サンドイッチである。パンがふっくらしていておいしい。

  ものの数分で平らげて、公園を去った。
